% Paper 4: Monetary Predator-Prey Dynamics: Enforcement Gridlock and
%          Neutral Settlement Survival
% Build: xelatex BGT-PAPER-4.tex (run twice for cross-references)

\documentclass[12pt]{article}

% --- Math (must load before unicode-math) ---
\usepackage{mathtools}  % loads amsmath automatically

% --- Fonts and encoding ---
\usepackage{fontspec}
\setmainfont{Latin Modern Roman}
\usepackage{unicode-math}
\setmathfont{Latin Modern Math}

% --- Page layout ---
\usepackage[margin=1in]{geometry}
\usepackage{setspace}
\onehalfspacing

% --- Theorem environments ---
\usepackage{amsthm}
\newtheorem{theorem}{Theorem}
\newtheorem{proposition}[theorem]{Proposition}
\newtheorem{lemma}[theorem]{Lemma}
\newtheorem{corollary}[theorem]{Corollary}
\theoremstyle{definition}
\newtheorem{definition}{Definition}
\newtheorem{assumption}{Assumption}
\newtheorem{condition}{Condition}
\theoremstyle{remark}
\newtheorem*{remark}{Remark}

% --- Tables ---
\usepackage{booktabs}
\usepackage{longtable}
\usepackage{array}
\usepackage{tabularx}
\usepackage{calc}

% --- Misc ---
\usepackage{enumitem}
\usepackage{hyperref}
\hypersetup{colorlinks=true, linkcolor=black, citecolor=black, urlcolor=blue}
\usepackage{microtype}

% --- Suppress section numbering (manual numbers in headings) ---
\setcounter{secnumdepth}{0}

\begin{document}

% ============================================================
% TITLE BLOCK
% ============================================================

\begin{center}
{\LARGE\bfseries Monetary Predator-Prey Dynamics:\\Enforcement Gridlock and Neutral Settlement Survival}

\bigskip

{\large Sean Hash}\\
bitcoingametheory.com\\
sean@bitcoingametheory.com

\bigskip

February 2026

\medskip

{\small
\textbf{JEL Codes:} C72 (Noncooperative Games), D74 (Conflict; Conflict Resolution), E42 (Monetary Systems), F33 (International Monetary Arrangements), G15 (International Financial Markets)\\[4pt]
\textbf{Keywords:} Bitcoin, Lotka--Volterra, enforcement gridlock, predator-prey dynamics, neutral settlement, coordination failure, coalition stability, game theory
}
\end{center}

\bigskip

% ============================================================
% ABSTRACT
% ============================================================

\begin{abstract}
\noindent
This paper formalizes the survivability question for neutral settlement:
\emph{why can enforcement actors not coordinate to suppress it?} We model
enforcement actors as ``predators'' (coordination taxers) and neutral
settlement rail usage as ``prey'' within a multi-predator Lotka--Volterra
framework. When inter-predator competition coefficients are strictly
positive---a condition guaranteed by geopolitical multipolarity---the
interior equilibrium ensures prey survival ($x^{*} > 0$). We prove that
no stable suppression coalition exists: for any proposed coalition, at
least one member gains by defecting (Theorem~\ref{thm:gridlock}).
Further, we establish that each enforcer's dominant strategy is to
preserve neutral rail access as a hedge against rival enforcers' monetary
weaponization (Lemma~\ref{lem:hedging}). Finally, we introduce duration
fragility inequalities (Proposition~\ref{prop:duration}) showing that
equity valuations suffer terminal value collapse under AI-driven moat
erosion, strengthening the marginal case for scarcity-based settlement.
These results complete the causal chain from Hash~(2026a): the
coordination failure in the exit game is not merely possible but
\emph{structurally guaranteed} by inter-predator competition.
\end{abstract}

\newpage

% ============================================================
\section{1. Introduction}\label{sec:intro}
% ============================================================

The Bitcoin Game Theory framework establishes that exit to neutral
settlement is self-reinforcing (Hash, 2026a, Theorem~1), that
coordination to stay fails (Theorem~2), and that the resulting
equilibrium is absorbing (Theorem~3). The seven-property elimination
(Hash, 2026b) shows Bitcoin is the unique survivor among candidate
assets. The trust gradient (Hash, 2026c) extends the result to
autonomous economic agents.

A critical gap remains: \emph{why} does coordination to stay fail?  The
prior result (Hash, 2026a, Theorem~2) establishes that no enforcement
mechanism exists to prevent exit, but does not explain the structural
dynamics that guarantee this condition persists. This paper fills that
gap.

We observe that enforcement actors---sovereigns imposing capital
controls, corporations gatekeeping financial access, intermediaries
extracting compliance rents---are not a monolithic bloc cooperating
against Bitcoin. They are \emph{competitors}, each taxing coordination
in ways that interfere with other enforcers' interests. This competitive
structure maps precisely onto multi-predator Lotka--Volterra population
dynamics (Volterra, 1928; May, 1973), where inter-species competition
among predators guarantees prey survival at the interior equilibrium.

\subsection{1.1 Contributions}

\begin{enumerate}[label=(\roman*)]
  \item \textbf{Actor Taxonomy} (Section~2): We classify monetary actors
    into coordination taxers (CT$_1$--CT$_3$) and exit-valve
    participants (EV$_1$--EV$_3$), identifying dual-role actors who
    simultaneously enforce and exit.
  \item \textbf{Lotka--Volterra Formalization} (Section~3): We model the
    system as a multi-predator ecology and derive existence conditions
    for the interior equilibrium.
  \item \textbf{Gridlock Wedge Theorem} (Section~4): We prove that
    inter-predator competition permanently prevents coordinated
    suppression (Theorem~\ref{thm:gridlock}).
  \item \textbf{Predator Hedging Lemma} (Section~5): We establish the
    dominant strategy cascade in which each enforcer preserves the
    neutral rail (Lemma~\ref{lem:hedging}).
  \item \textbf{Duration Fragility} (Section~6): We prove that
    AI-driven moat erosion compresses equity terminal values
    (Proposition~\ref{prop:duration}).
  \item \textbf{Falsification} (Section~8): We specify condition~F7
    under which these results fail.
\end{enumerate}

% ============================================================
\section{2. Actor Taxonomy}\label{sec:taxonomy}
% ============================================================

\subsection{2.1 Coordination Taxers}

\begin{definition}[Coordination Taxer]\label{def:ct}
A \emph{coordination taxer} (CT) is an enforcement actor that extracts
value by controlling or taxing monetary flows. We identify three tiers:
\end{definition}

\begin{table}[h]
\centering
\begin{tabular}{@{}clll@{}}
\toprule
ID & Tier & Examples & Mechanism \\
\midrule
CT$_1$ & Sovereign & Governments, central banks, & Monetary policy, capital \\
       &           & sanctions agencies           & controls, legal coercion \\[3pt]
CT$_2$ & Corporate & Banks, asset managers,       & Intermediation fees, \\
       &           & payment processors           & compliance gatekeeping \\[3pt]
CT$_3$ & Intermediary & Prime brokers, mining     & Collateral yield, pool \\
       &              & pools, payment routers    & fees, routing fees \\
\bottomrule
\end{tabular}
\caption{Coordination taxer classification.}
\label{tab:ct}
\end{table}

\subsection{2.2 Exit-Valve Participants}

\begin{definition}[Exit-Valve Participant]\label{def:ev}
An \emph{exit-valve participant} (EV) is an actor who benefits from
access to a neutral settlement rail outside the enforcement perimeter.
\end{definition}

\begin{table}[h]
\centering
\begin{tabular}{@{}clll@{}}
\toprule
ID & Tier & Examples & Mechanism \\
\midrule
EV$_1$ & Retail & Savers, remittance users, & Inflation escape, bank \\
       &        & unbanked populations      & freeze bypass \\[3pt]
EV$_2$ & Tier-2 Sovereign & Rival powers, & Dollar alternative, \\
       &                  & sanctioned states & sanctions evasion \\[3pt]
EV$_3$ & Corporate & Energy producers, miners, & Energy monetization, \\
       &           & AI agents                & autonomous settlement \\
\bottomrule
\end{tabular}
\caption{Exit-valve participant classification.}
\label{tab:ev}
\end{table}

\subsection{2.3 Dual-Role Actors}

Many actors simultaneously tax coordination \emph{and} use the neutral
rail, creating the strategic tension that drives the Gridlock Wedge:

\begin{itemize}
  \item \textbf{United States} (CT$_1$ + EV$_2$): Enforces dollar
    hegemony via SWIFT sanctions while maintaining a strategic Bitcoin
    reserve as a hedge against de-dollarization.
  \item \textbf{China} (CT$_1$ + EV$_2$): Imposes capital controls and
    bans domestic crypto trading while using cryptocurrency for
    cross-border trade settlement.
  \item \textbf{Russia/Iran} (CT$_1$ + EV$_2$): Nominally restrict
    domestic use while employing Bitcoin for sanctions evasion in
    international trade.
  \item \textbf{Major banks} (CT$_2$ + EV$_3$): Gatekeep fiat access
    while building custody and trading infrastructure to capture
    Bitcoin-denominated revenue.
\end{itemize}

\begin{remark}
In informal discussion, coordination taxers are sometimes called
``cats'' and exit-valve participants ``mice,'' referencing the familiar
cat-and-mouse enforcement dynamic. We use formal terms throughout.
\end{remark}

% ============================================================
\section{3. Model: Multi-Predator Lotka--Volterra System}\label{sec:model}
% ============================================================

\subsection{3.1 System Equations}

Let $x(t) \in \mathbb{R}_{\geq 0}$ denote neutral rail usage (a proxy
for prey population) and $y_k(t) \in \mathbb{R}_{\geq 0}$ the
enforcement intensity of actor~$k \in \{1, \ldots, n\}$.

\begin{definition}[Multi-Predator Enforcement System]\label{def:system}
The enforcement dynamics are governed by:
\begin{align}
  \frac{dx}{dt} &= r\,x\!\left(1 - \frac{x}{K}\right)
    - \sum_{k=1}^{n} a_k \, x \, y_k \,,
    \label{eq:prey} \\[6pt]
  \frac{dy_k}{dt} &= b_k \, x \, y_k - d_k \, y_k
    - \sum_{\substack{j=1 \\ j \neq k}}^{n}
      \varepsilon_{jk} \, y_j \, y_k \,,
    \qquad k = 1, \ldots, n \,,
    \label{eq:predator}
\end{align}
where the parameters are defined in Table~\ref{tab:params}.
\end{definition}

\begin{table}[h]
\centering
\begin{tabular}{@{}cl@{}}
\toprule
Symbol & Definition \\
\midrule
$x(t)$            & Neutral rail usage (prey population proxy) \\
$y_k(t)$          & Enforcement intensity of actor~$k$ \\
$r > 0$           & Intrinsic growth rate of neutral rail adoption \\
$K > 0$           & Carrying capacity (maximum adoption) \\
$a_k > 0$         & Suppression efficiency of enforcer~$k$ \\
$b_k > 0$         & Benefit enforcer~$k$ extracts from enforcement \\
$d_k > 0$         & Cost/decay rate of enforcement effort \\
$\varepsilon_{jk} > 0$ & Inter-predator competition coefficient between $j$ and~$k$ \\
\bottomrule
\end{tabular}
\caption{Model parameters.}
\label{tab:params}
\end{table}

\subsection{3.2 Inter-Predator Competition}

The critical structural parameter is $\varepsilon_{jk}$: the degree to
which enforcer~$j$'s actions interfere with enforcer~$k$'s enforcement
capacity.

\begin{assumption}[Multipolarity]\label{asm:multi}
There exist at least three tier-1 enforcement actors ($n \geq 3$), and
inter-predator competition is strictly positive:
$\varepsilon_{jk} > 0$ for all $j \neq k$.
\end{assumption}

\begin{remark}
Assumption~\ref{asm:multi} is not a modeling convenience---it is an
empirical regularity guaranteed by geopolitical multipolarity (AX1/W1 in
Hash, 2026a). Concrete channels include: US sanctions pushing sanctioned
parties toward Bitcoin (undermining China's capital controls); EU
regulation pushing capital to permissive jurisdictions (undermining
enforcement coherence); bank compliance requirements creating demand for
non-custodial alternatives (expanding the enforcement perimeter).
\end{remark}

\subsection{3.3 Interior Equilibrium}

Setting \eqref{eq:prey}--\eqref{eq:predator} to zero and solving for
the interior fixed point $(x^{*}, y_1^{*}, \ldots, y_n^{*})$ with all
components strictly positive, we obtain:

\begin{equation}\label{eq:xstar}
  x^{*} = K \left( 1 - \frac{1}{r} \sum_{k=1}^{n} a_k \, y_k^{*}
  \right) > 0 \,,
\end{equation}

where the $y_k^{*}$ satisfy the coupled system
\begin{equation}\label{eq:ystar}
  b_k \, x^{*} - d_k - \sum_{\substack{j=1 \\ j \neq k}}^{n}
    \varepsilon_{jk} \, y_j^{*} = 0 \,,
    \qquad k = 1, \ldots, n \,.
\end{equation}

The positivity of $x^{*}$ in \eqref{eq:xstar} holds because
inter-predator friction bounds each $y_k^{*}$ from above: the
competition terms $\varepsilon_{jk} \, y_j^{*}$ in
\eqref{eq:ystar} prevent any single enforcer from achieving the
suppression threshold at which $\sum_k a_k y_k^{*} / r \geq 1$.

% ============================================================
\section{4. The Gridlock Wedge Theorem}\label{sec:gridlock}
% ============================================================

\begin{theorem}[Gridlock Wedge]\label{thm:gridlock}
Under Assumption~\ref{asm:multi} ($n \geq 3$,
$\varepsilon_{jk} > 0$ for all $j \neq k$), the following hold:
\begin{enumerate}[label=\textup{(G\arabic*)}]
  \item\label{G1} Multiple enforcers exist at all observed~$t$.
  \item\label{G2} Inter-predator competition is strictly positive.
  \item\label{G3} The neutral settlement rail survives at a positive
    interior equilibrium: $x^{*} > 0$.
  \item\label{G4} Each enforcer strictly prefers that rivals not
    suppress.
  \item\label{G5} Given that rivals preserve the rail, each enforcer's
    best response is to also preserve.
  \item\label{G6} No stable suppression coalition exists.
\end{enumerate}
\end{theorem}

\begin{proof}
\ref{G1} is an empirical observation: at least three tier-1 enforcement
actors (US, EU, China) operate simultaneously at all observed time
periods.

\ref{G2} follows from \ref{G1} and Assumption~\ref{asm:multi}: every
enforcer pair exhibits strictly positive competition
($\varepsilon_{jk} > 0$) through the channels described in
Section~3.2.

For \ref{G3}, rearrange each predator's equilibrium
condition~\eqref{eq:ystar}:
\begin{equation}\label{eq:survival}
  b_k \, x^{*} = d_k + \sum_{\substack{j=1 \\ j \neq k}}^{n}
    \varepsilon_{jk} \, y_j^{*} \,.
\end{equation}
Since $d_k > 0$ and every term on the right-hand side is non-negative,
\[
  x^{*} = \frac{d_k + \sum_{j \neq k}
    \varepsilon_{jk} \, y_j^{*}}{b_k}
    \;\geq\; \frac{d_k}{b_k} \;>\; 0 \,.
\]
This is the standard survival result in multi-predator
Lotka--Volterra theory: each predator's equilibrium condition
\emph{individually} guarantees $x^{*} > 0$, because enforcement
profitability requires a positive prey population.  Inter-predator
competition strengthens this bound: the $\varepsilon_{jk} \, y_j^{*}$
terms in~\eqref{eq:survival} increase $x^{*}$ above the
single-predator baseline $d_k/b_k$, as competition diverts enforcement
resources from prey suppression to inter-predator rivalry.

For \ref{G4}, each enforcer~$k$'s per-period extraction is
$b_k \, x \, y_k$, which is strictly increasing in~$x$.  If rivals
intensify suppression (raise $y_j$ for $j \neq k$), the prey
equation~\eqref{eq:prey} yields a lower equilibrium~$x$, reducing
$k$'s extraction.  Conversely, if rivals \emph{reduce} enforcement,
$x$ rises and $k$'s extraction improves.  Thus each enforcer strictly
prefers that rivals maintain lower enforcement intensity:
\[
  \frac{\partial}{\partial y_j}\!\bigl(b_k \, x^{*}(y_j)\bigr)
  < 0 \quad \text{for all } j \neq k \,,
\]
where $x^{*}(y_j)$ denotes the equilibrium prey level as a function of
rival enforcement.  No surplus-splitting mechanism is required; the
preference follows directly from the competitive structure of the
multi-predator system.

\ref{G5} follows from game-theoretic dominance analysis. Fix the
strategies of all enforcers $j \neq k$ at ``preserve.'' If $k$ bans,
$k$ loses neutral rail access (hedge value against rivals'
weaponization) while competitors maintain access. If $k$ preserves, $k$
retains both hedge value and compliance revenue from supervised activity.
Since the hedge value is strictly positive under multipolarity (rival
weaponization is a non-degenerate threat), preservation weakly dominates
banning. This generates a dominant strategy cascade: once any enforcer
preserves, all others' best response is to preserve
(Lemma~\ref{lem:hedging}).

For \ref{G6}, suppose a suppression coalition
$C \subseteq \{1, \ldots, n\}$ forms. Consider any member~$k \in C$.
If $k$ deviates to ``preserve'' while $C \setminus \{k\}$ continues
suppressing:
\begin{enumerate}[label=(\alph*)]
  \item $k$ captures fleeing capital from the suppressed jurisdictions.
  \item $k$ retains neutral rail access as a hedge against
    $(C \setminus \{k\})$'s residual enforcement power.
  \item Rivals in $C \setminus \{k\}$ bear higher per-capita
    enforcement costs due to $k$'s defection.
\end{enumerate}
The deviation payoff strictly exceeds the coalition payoff, so no Nash
stable coalition exists.
\end{proof}

% ============================================================
\section{5. The Dominant Strategy Cascade}\label{sec:hedging}
% ============================================================

\begin{lemma}[Predator Hedging]\label{lem:hedging}
Let $\mathcal{G}$ be the simultaneous-move game in which each enforcer
$k \in \{1, \ldots, n\}$ chooses $s_k \in \{\text{ban}, \text{preserve}\}$.
Under Assumption~\ref{asm:multi}, ``preserve'' weakly dominates ``ban''
for every~$k$, regardless of rivals' strategies. Consequently,
suppression is never unanimous.
\end{lemma}

\begin{proof}
Define enforcer~$k$'s payoff function as
\begin{equation}\label{eq:payoff}
  U_k(s_k, s_{-k}) = \underbrace{R_k(s_k, s_{-k})}_{\text{enforcement revenue}}
    + \underbrace{H_k(s_k, s_{-k})}_{\text{hedge value}}
    - \underbrace{C_k(s_k)}_{\text{enforcement cost}} \,,
\end{equation}
where:
\begin{itemize}
  \item $R_k$ is the revenue from enforcement activity (compliance fees,
    taxation of supervised Bitcoin activity).
  \item $H_k$ is the option value of maintaining neutral rail access as
    a hedge against rival monetary weaponization (sanctions, SWIFT
    exclusion, capital freezes).
  \item $C_k$ is the direct cost of the chosen strategy.
\end{itemize}

\medskip\noindent
\textbf{Case 1:} Rival strategies $s_{-k}$ include at least one
``preserve.''

If some rival preserves, the neutral rail persists. If $k$~bans,
$H_k = 0$ (loses hedge). If $k$~preserves, $H_k > 0$ and $R_k \geq 0$
(can tax supervised activity). Since $C_k(\text{preserve}) \leq
C_k(\text{ban})$ (suppression requires active enforcement),
$U_k(\text{preserve}) > U_k(\text{ban})$.

\medskip\noindent
\textbf{Case 2:} All rivals choose ``ban'' ($s_j = \text{ban}$ for all
$j \neq k$).

If all rivals ban and $k$ also bans, the rail is fully suppressed:
$R_k = 0$, $H_k = 0$, but $C_k(\text{ban}) > 0$. If $k$ deviates to
``preserve,'' $k$ captures all fleeing capital from rival jurisdictions
($R_k \gg 0$) and retains the hedge ($H_k > 0$). Thus
$U_k(\text{preserve}) \gg U_k(\text{ban})$.

\medskip
In both cases, ``preserve'' yields a payoff at least as large as
``ban.'' Since Case~2 yields strict inequality, ``preserve'' is the
unique best response for at least one~$k$, which by induction
(Case~1) propagates to all enforcers.
\end{proof}

\subsection{5.1 Empirical Validation}

The predator hedging mechanism is observable in practice:

\begin{itemize}
  \item \textbf{Russia (2022--present):} Despite nominally restricting
    cryptocurrency domestically, Russia authorized Bitcoin for
    international settlements to circumvent Western
    sanctions---exhibiting dual-role behavior (CT$_1$ + EV$_2$).
  \item \textbf{El~Salvador (2021--present):} Small sovereign adopts
    Bitcoin as legal tender, demonstrating EV$_2$ first-mover behavior;
    IMF pressure fails to reverse the policy.
  \item \textbf{United States (2024--present):} Establishes strategic
    Bitcoin reserve while maintaining SEC enforcement---CT$_1$ + EV$_2$
    simultaneously.
\end{itemize}

% ============================================================
\section{6. Duration Fragility}\label{sec:duration}
% ============================================================

We now establish that the principal alternative store of value---equity
in AI companies---faces structural vulnerabilities that strengthen the
case for scarcity-based settlement.

\subsection{6.1 Moat Half-Life Under AI Acceleration}

\begin{definition}[Moat Half-Life]\label{def:moat}
Let $H(t)$ denote the competitive moat half-life of a technology firm
at time~$t$. Under AI-driven acceleration:
\begin{equation}\label{eq:moat}
  H(t) = H_0 \, e^{-\alpha t} \,,
\end{equation}
where $H_0 > 0$ is the initial moat half-life and $\alpha > 0$ captures
the rate of AI-driven competitive erosion.
\end{definition}

Historical calibration suggests $H_0 \approx 15$ years for pre-AI
technology moats (e.g., Microsoft's OS dominance persisted for
approximately~15 years before commoditization pressure materialized). Under
AI acceleration with $\alpha \approx 0.20$, projected moat half-lives
compress to $H \approx 3$--$5$ years.

\subsection{6.2 Multiple Compression}

\begin{proposition}[Duration Fragility]\label{prop:duration}
Let $E_t$ denote a firm's current earnings, $r > 0$ the risk-free
discount rate, and $\alpha \geq 0$ the moat erosion rate. The
earnings-based terminal value satisfies:
\begin{equation}\label{eq:pe}
  V_t = \frac{E_t}{r + \alpha} \,,
\end{equation}
which is strictly decreasing in $\alpha$. In particular:
\begin{enumerate}[label=(\alph*)]
  \item At $\alpha = 0$ (no moat erosion): $P/E = 1/r$.
  \item At $\alpha = r$ (moat erodes at discount rate): $P/E = 1/(2r)$,
    a $50\%$ compression.
  \item As $\alpha \to \infty$ (instantaneous commoditization):
    $V_t \to 0$.
\end{enumerate}
\end{proposition}

\begin{proof}
The terminal value of a perpetual earnings stream $E_t$ subject to
exponential decay at rate~$\alpha$ is
\[
  V_t = \int_0^{\infty} E_t \, e^{-\alpha s} \, e^{-rs} \, ds
      = E_t \int_0^{\infty} e^{-(r+\alpha)s} \, ds
      = \frac{E_t}{r + \alpha} \,.
\]
Statements (a)--(c) follow by substitution.
\end{proof}

\begin{remark}[The 100$\times$ Compression Scenario]
A firm earning \$10/share with $r = 0.10$ and $\alpha = 0$ trades at
$P/E = 10$ ($\$100$/share). With $\alpha = 0.20$ (moat half-life
$\approx 3.5$~years): $P/E = 3.3$ ($\$33$/share)---a $67\%$ haircut
with \emph{identical} current earnings. At $\alpha = 0.90$: $P/E = 1.0$
($\$10$/share)---a $90\%$ haircut. The fiber-optic precedent
(1997--2010: margins from $60\%$ to $15\%$) demonstrates this pattern at
infrastructure scale.
\end{remark}

\subsection{6.3 Implication for Settlement Asset Selection}

Bitcoin has no earnings and therefore no duration fragility. Its value
derives entirely from scarcity (21M fixed supply) and network effects
(Metcalfe-type scaling)---properties immune to competitive moat erosion.
For long-duration wealth preservation across generational time horizons,
the absence of duration fragility is a structural advantage over any
earnings-dependent asset class.

% ============================================================
\section{7. Empirical Evidence}\label{sec:evidence}
% ============================================================

\begin{table}[h]
\centering
\small
\begin{tabular}{@{}p{5cm}cp{3.2cm}@{}}
\toprule
Observation & Classification & Supports \\
\midrule
Russia authorizes crypto for international trade (2024)
  & CT$_1$ + EV$_2$
  & Lemma~\ref{lem:hedging}, \ref{G5} \\[3pt]
Iran uses Bitcoin mining for energy monetization (2020--)
  & CT$_1$ + EV$_3$
  & \ref{G4} \\[3pt]
El Salvador accumulates 6{,}000+ BTC (2021--2025)
  & EV$_2$
  & E4, K1 \\[3pt]
US establishes strategic Bitcoin reserve (2024--)
  & CT$_1$ + EV$_2$
  & Lemma~\ref{lem:hedging}, \ref{G5} \\[3pt]
China bans crypto trading but mines hash (2021--)
  & CT$_1$ + EV$_2$/EV$_3$
  & \ref{G2}, \ref{G4} \\[3pt]
NVIDIA P/E: $70\times \to 30\times$ (2023--2026)
  & ---
  & Prop.~\ref{prop:duration} \\[3pt]
Fiber-optic margins: $60\% \to 15\%$ (1997--2010)
  & ---
  & Historical precedent \\[3pt]
Singapore issues 13 DPT licenses (2024)
  & EV$_2$
  & \ref{G2}, \ref{G4} \\
\bottomrule
\end{tabular}
\caption{Empirical observations classified by actor taxonomy.}
\label{tab:evidence}
\end{table}

% ============================================================
\section{8. Falsification}\label{sec:falsification}
% ============================================================

\begin{condition}[F7: Gridlock Closes]\label{cond:F7}
The results of this paper are falsified if:
\begin{quote}
\emph{Synchronized global suppression eliminates all enforcement gaps
\textbf{and} permanent tier-1 capability lockout prevents
re-emergence.}
\end{quote}
\end{condition}

Specifically, F7 requires:
\begin{enumerate}[label=(\roman*)]
  \item All tier-1 enforcers (US, EU, China at minimum) simultaneously
    and permanently suppress Bitcoin access.
  \item No tier-2 sovereign defects to capture fleeing capital.
  \item Technical capability to run nodes and mine is permanently
    eliminated (not merely pushed underground).
  \item The above conditions persist indefinitely (not merely during a
    crisis period).
\end{enumerate}

F7 subsumes and strengthens F4 (stable cartel, Hash 2026a): it requires
not just a cartel but synchronized action with permanent capability
destruction. Under Assumption~\ref{asm:multi}, the inter-predator
competition coefficients $\varepsilon_{jk}$ would need to collapse to
zero---meaning geopolitical competition itself would need to end.

% ============================================================
\section{9. Limitations}\label{sec:limitations}
% ============================================================

\begin{enumerate}[label=(\roman*)]
  \item \textbf{Parameter estimation:} The Lotka--Volterra coefficients
    ($a_k$, $b_k$, $d_k$, $\varepsilon_{jk}$) are modeled
    qualitatively. Empirical calibration would strengthen the results
    but is not required for the qualitative survival conclusion
    ($x^{*} > 0$).
  \item \textbf{Regime changes:} The model assumes multipolarity
    persists. If W1 (Open World) transitions to W2 (Closed World---a
    single global hegemon), the inter-predator competition coefficients
    collapse to zero and F7 becomes possible.
  \item \textbf{Duration fragility calibration:} The moat erosion rate
    $\alpha$ is estimated from limited historical precedent. If AI
    creates genuinely permanent moats (unprecedented in technological
    history), the equity compression argument weakens.
  \item \textbf{Biological analogy limits:} Monetary enforcement
    dynamics differ from biological predator-prey systems in that
    actors are strategic (game-theoretic) rather than mechanistic. The
    Lotka--Volterra structure captures the competitive dynamics but
    understates the full strategic complexity. A richer formulation
    using differential games (Isaacs, 1965) could refine the
    quantitative predictions without altering the qualitative survival
    result.
\end{enumerate}

% ============================================================
\section{10. Conclusion}\label{sec:conclusion}
% ============================================================

The enforcement coordination problem is not merely a prisoner's
dilemma---it is a structurally guaranteed gridlock arising from
inter-predator competition. Competing enforcement actors cannot suppress
the neutral settlement rail because each preserves it as a hedge against
rivals. This result completes the causal chain:

\begin{enumerate}
  \item Exit is self-reinforcing (Hash, 2026a, Theorem~1).
  \item Coordination to stay fails \emph{because enforcement actors are
    in permanent gridlock} (Theorem~\ref{thm:gridlock}).
  \item The resulting equilibrium is absorbing (Hash, 2026a, Theorem~3).
  \item Bitcoin is the unique survivor of seven-property elimination
    (Hash, 2026b).
  \item The result extends to autonomous agents at zero trust (Hash,
    2026c).
\end{enumerate}

The correct framing for Bitcoin's survivability is not eradication risk
but \emph{containment}. Enforcement actors will regulate, tax, and
monitor Bitcoin activity at the application layer---but they cannot and
will not eliminate the settlement layer, because doing so would surrender
their hedge against rival enforcers. The neutral rail survives not
despite enforcement but \emph{because of} the competitive structure of
enforcement itself.

% ============================================================
\section{References}\label{sec:references}
% ============================================================

\begin{flushleft}

Chainalysis (2024). The 2024 Crypto Crime Report.

\medskip

Hash (2026a). Bitcoin exit dominance in monetary coordination games.
Working paper, bitcoingametheory.com.

\medskip

Hash (2026b). Bitcoin as unique neutral settlement: A seven-property
elimination. Working paper, bitcoingametheory.com.

\medskip

Hash (2026c). Settlement at zero trust: Bitcoin and autonomous economic
agents. Working paper, bitcoingametheory.com.

\medskip

Isaacs, R. (1965). \emph{Differential Games: A Mathematical Theory with
Applications to Warfare and Pursuit, Control and Optimization.} John
Wiley \& Sons.

\medskip

May, R. M. (1973). \emph{Stability and Complexity in Model Ecosystems.}
Princeton University Press.

\medskip

Volterra, V. (1928). Variations and fluctuations of the number of
individuals in animal species living together. \emph{ICES Journal of
Marine Science}, 3(1), 3--51.
\url{https://doi.org/10.1093/icesjms/3.1.3}

\end{flushleft}

% ============================================================
\section{Author's Note}\label{sec:author}
% ============================================================

``Sean Hash'' is a pen name. The author is a single human individual
writing under a pseudonym. The framework is designed to be evaluated on
its logical structure---every claim derives from four empirical axioms,
and every claim specifies falsification conditions. The author's identity
is irrelevant to whether the axioms hold.

The author acknowledges modest priors: a preference for multipolar worlds
over global monopoly, and for individual economic sovereignty over total
state control of capital. These do not predetermine the framework's
conclusions---the game-theoretic derivation works for purely selfish
actors regardless of values.

The author also acknowledges a real moral cost: a neutral settlement layer
will process transactions for harmful actors. This cost is real, but
settlement is not acceptance---the acceptance game produces a natural
enforcement equilibrium through rational counterparty decisions, without
requiring protocol-level censorship.

For full disclosure on authorship, moral position, and the
settlement-acceptance distinction, see the Author's Note
(bitcoingametheory.com/rfc/BGT-AUTHOR.txt).

\section{Disclosure}\label{sec:disclosure}
% ============================================================
% APPENDIX
% ============================================================

\appendix
\section{Appendix: Notation}\label{app:notation}

\begin{table}[h]
\centering
\begin{tabular}{@{}cl@{}}
\toprule
Symbol & Definition \\
\midrule
$x(t)$                 & Neutral rail usage (prey population proxy) \\
$y_k(t)$               & Enforcement intensity of actor~$k$ \\
$n$                    & Number of enforcement actors \\
$r$                    & Intrinsic growth rate of neutral rail adoption \\
$K$                    & Carrying capacity (maximum adoption) \\
$a_k$                  & Suppression efficiency of enforcer~$k$ \\
$b_k$                  & Benefit enforcer~$k$ extracts from enforcement \\
$d_k$                  & Cost/decay rate of enforcement effort for~$k$ \\
$\varepsilon_{jk}$     & Inter-predator competition coefficient ($j \neq k$) \\
$x^{*}, y_k^{*}$      & Interior equilibrium values \\
$H(t)$                 & Competitive moat half-life at time~$t$ \\
$H_0$                  & Initial moat half-life (pre-AI) \\
$\alpha$               & Rate of AI-driven moat erosion \\
$E_t$                  & Current earnings of a firm \\
$V_t$                  & Terminal value of earnings stream \\
$U_k(s_k, s_{-k})$    & Enforcer~$k$'s payoff under strategy profile \\
$R_k, H_k, C_k$       & Revenue, hedge value, and cost components \\
CT$_1$--CT$_3$         & Coordination taxer tiers \\
EV$_1$--EV$_3$         & Exit-valve participant tiers \\
\bottomrule
\end{tabular}
\caption{Summary of notation.}
\label{tab:notation}
\end{table}

\end{document}
